%%%%%
%
\section{Correction de sauts dans le temps acquis par le Merlin IV}

\itemstyle{Op\'erateur~:} {\sf correc\_time\_rec}

%%%%%OBJET%%%%
\itemstyle{Objet~:}
Correction des sauts anormaux dans le temps acquis par le
syst\`eme d'acquisition du Merlin IV. 

%%%%%ARGUMENTS%%%%
\bigskip
%%%%%ENTREES%%%%
\begin{tabarg}
\multicolumn{3}{l}{\titretab {Entr\'ees}}\\
\hline
$Time$ & Vecteur, {\tt double} & Temps ($s$)\\
$f_e$ & Coeff., {\tt float} & Fr\'equence d'\'echantillonnage ($Hz$)\\
\end{tabarg}

\medskip
%%%%%SORTIES%%%%
\begin{tabarg}
\multicolumn{3}{l}{\titretab {Sorties}}\\
\hline
$Time_c$ & Vecteur, {\tt double} & Temps corrig\'e ($s$)\\
\end{tabarg}


%%%%%FORMULE%%%%
\itemstyle{Formule}

%%correc_time_rec
Correction des sauts de $Time$.


%%%% COMMENTAIRE %%%%
\itemstyle{Commentaire~:}

En principe, le temps cro\^{\i}t r\'eguli\`erement
de $1/f_e \pm \varepsilon $ \`a chaque instant.
Mais, il peut se produire des "sauts" ou des "trous".
Un "saut" est constitu\'e par N valeurs 
cons\'ecutives "incorrectes". 
Un "trou" est l'absence de N valeurs 
cons\'ecutives dans la s\'erie. \\
L'algorithme remplace les valeurs "incorrectes" des
"sauts" et ne touche pas les "trous".\\
$\varepsilon$ est l'incertitude sur le temps du
syst\`eme d'acquisition. Elle doit \^etre ici
inf\'erieure \`a $6.10^{-4} s$ ( sp\'ecifique au
Merlin IV) pour que la valeur soit consid\'er\'ee
correcte.
Un "saut" ne peut pas durer plus de 20 valeurs
cons\'ecutives sinon il est consid\'er\'e 
comme un "trou".


%%%%%AUTEUR%%%%
\itemstyle{Auteur~:}
CNRM/GMEI/TRAMM

%%%%%REFERENCE%%%%
\itemstyle{R\'ef\'erences~:}
